\section{Discussion}
\label{sec:discussion}

\subsection{Mechanism of sextic promotion}

The central feature of the resonance is not an unusually large sixth-order term by itself. It is the suppression of the quartic term in a distinguished angular sector.

Along the axes, the quartic Hamiltonian remains nondegenerate and yields the ordinary law
\[
|\mcJ_H|\asymp\delta^2.
\]

Near the diagonals, the coefficient of the quartic radial term is proportional to the small quantity
\[
\sigma=4-\rho^2.
\]
The sixth-order term therefore enters at leading order and generates the intermediate law
\[
|\mcJ_E|\asymp\delta^{3/2}.
\]

This is the sense in which the sextic term is dynamically promoted.

\subsection{Why the quartic balance is meaningful}

The quartic balance is not imposed by an arbitrary anisotropic rescaling. The detuning fixes the linear symplectic scale, and the reversing involutions fix the remaining rotational freedom.

An independently defined quartic balancing scale agrees numerically with the detuning scale to very high accuracy. The additional orbital reconstruction of \(A\) and \(\rho\) provides a second test of the same conclusion.

\subsection{Intermediate scaling and crossover}

For fixed
\[
\sigma<0,
\]
the \(3/2\) law is not the ultimate \(\delta\to0\) asymptotic. The correct interpretation is
\[
\deltax\ll\delta\ll1
\]
for the sextic-dominated regime, followed by a crossover to a finite-amplitude plateau.

The normalized scalar model provides a compact description of that transition.

\subsection{Residual elliptic cycles}

The negative sign of the quartic diagonal coefficient leads, in the sixth-order truncation, to nonzero elliptic equilibria at resonance. The exact second return numerically exhibits the corresponding residual period-two cycles.

Their existence, multiplicity, elliptic stability, action scale, and Floquet scale together provide a restrictive test of the sixth-order model.

\subsection{Observable accessibility of the normal form}

The leading axial action determines \(A\), while the leading physical second-return Floquet exponent determines \(\rho\).  These two limits are robust because the first detuning-dependent corrections enter only one order later, and the Richardson extrapolations converge toward the critical-jet values.

The numerical audit also clarifies what cannot be inferred from the same axial data.  At subleading order, the critical sextic Hamiltonian is mixed with detuning-dependent quartic terms and with quadratic corrections to the linear detuning; see \cref{sec:axial-contamination,app:axial-contamination}.  We therefore do not interpret the axial \(\delta^3\) action coefficients or \(\delta^2\) Floquet corrections as direct measurements of \(d,e,f,g\).

The diagonal sextic combination \(G\) is instead tested by the near-diagonal elliptic branch, for which the promoted sextic term enters asymptotically ahead of the first detuning-dependent quartic correction.  The observable hierarchy used in the final analysis is therefore
\[
\text{leading axial action}\longrightarrow A,
\]
\[
\text{leading axial Floquet exponent}\longrightarrow\rho,
\]
and
\[
\text{near-diagonal action--Floquet crossover}\longrightarrow G.
\]

The return jet and the periodic-orbit geometry remain complementary descriptions of the same resonant germ, but the comparison is made only at asymptotic orders where the relevant coefficient is identifiable.

\subsection{Model-dependent and coefficient-independent features}

The values of
\[
\mu_\star,A,B,d,e,f,g
\]
are specific to the present Carnot geodesic flow and the fixed reversible normal-form chart.

By contrast, once the scalar reduced Hamiltonian has the form
\[
-\delta I+aI^2+GI^3,
\qquad
a<0,\ G>0,
\]
the normalized relations
\[
Y=4X^3-3X^2
\]
and
\[
Z=X\sqrt{2X-1}
\]
are coefficient independent.

They should be regarded as leading scalar crossover relations, not as universal identities of arbitrary reversible \(1{:}2\) resonances.

\subsection{Scope}

The analysis is local in phase space. No conclusion about global integrability follows from the resonance alone.

Likewise, the exactly quartic-degenerate limit
\[
\sigma=0
\]
defines a different higher-codimension asymptotic problem. That limit is deliberately separated from the fixed-\(\sigma<0\) crossover studied here.
